\section{Conclusion}
\label{sec:conclusion}

We presented a geometric-aware approach for 6D object pose estimation that incorporates camera intrinsics and detection geometry to improve translation prediction. Our key contributions are:

\textbf{RGB Model (4-channel).} Using RGB+Mask input with ResNet-50 backbone, predicting rotation as quaternion and Z-depth only. X and Y are derived geometrically from the pinhole camera model. This achieves 41.2\% ADD-2cm accuracy with 31.5mm mean ADD error.

\textbf{RGB-D Model (5-channel).} A dual-stream fusion architecture that leverages depth sensor readings as a prior. The model predicts a small Z-offset from the sensor depth, rather than absolute depth. This achieves 84.9\% ADD-2cm accuracy with 13.5mm mean ADD error---more than 2$\times$ improvement over RGB-only.

\textbf{Key findings:}
\begin{itemize}
    \item Geometric translation (predict Z, derive X/Y) significantly reduces learning complexity
    \item Depth sensor provides strong prior: model learns only $\pm$10--30mm offset
    \item Homoscedastic uncertainty weighting eliminates manual loss weight tuning
    \item Four objects (cat, duck, eggbox, glue) achieve 100\% accuracy with RGB-D
    \item Iron and phone remain challenging due to reflective/texture-less surfaces
    \item Data augmentation (color jitter, bbox jitter) decreased accuracy, highlighting pose estimation's sensitivity to geometric constraints
\end{itemize}

\textbf{Quantitative Summary:}
\begin{itemize}
    \item ADD Error: RGB 31.5mm $\rightarrow$ RGB-D 13.5mm (-57.1\%)
    \item ACC@2cm: RGB 41.2\% $\rightarrow$ RGB-D 84.9\% (+106.1\%)
\end{itemize}

\textbf{Future Work.} Potential improvements include iterative refinement networks, attention-based feature fusion, and training on larger datasets with more diverse lighting and occlusion conditions. Additionally, extending to category-level pose estimation (novel object instances) would increase practical applicability.
